%% Chapitre 09 — Commandes, hooks et automatisations
%% RÈGLE : uniquement des commandes sémantiques bookforge. Aucun \chapter,
%% \textbf, \emph, \begin{lstlisting}, \vspace... brut.

\chapitre{Commandes, hooks et automatisations}

\introChapitre{
  Une fois les bases acquises, on industrialise ses gestes récurrents. Ce
  chapitre, daté car très lié à l'outillage du moment, présente les skills et
  commandes slash pour réutiliser des prompts, et les hooks pour automatiser des
  actions autour de l'agent. On transforme des manipulations répétées en
  raccourcis fiables.
}

\paragraphe{Au chapitre précédent, tu as appris à surveiller ce que coûte une
session et à couper le gaspillage. Le réflexe suivant est naturel : si tu
retapes le même prompt trois fois par jour, ou si tu relances à la main le
même formatage après chaque modification, c'est que le geste est mûr pour
l'automatisation. Ce chapitre outille exactement ça.}

\begin{avertissement}
Ce chapitre est \emphase{daté}. La syntaxe exacte des commandes, l'emplacement
des fichiers de configuration et le format des hooks évoluent vite d'une
version de Claude Code à l'autre. Les exemples ci-dessous illustrent
les \emphase{concepts}, pas une API figée. Avant de copier un extrait, vérifie
la documentation de ta version : lance \code{claude --version}, puis confronte
le format à la doc officielle courante. Si un champ ne fonctionne plus, c'est
le signe qu'il a changé de nom, pas que l'idée est fausse.
\end{avertissement}

\section{Skills et slash commands : encapsuler un prompt réutilisable}

\paragraphe{Dans Claude Code, un \terme{skill} est aujourd'hui la manière
conseillée d'encapsuler un prompt réutilisable. Tu le ranges dans
\code{.claude/skills/}, et il devient invocable comme une commande slash. Au
lieu de réécrire à chaque fois « relis ce diff, vérifie la gestion d'erreurs et
signale les cas limites oubliés », tu enregistres ce prompt une fois sous un nom
court et tu l'appelles avec \code{/relecture}.}

\paragraphe{L'intérêt n'est pas le gain de frappe. C'est la
\emphase{constance}. Un prompt qui traîne dans ta tête varie selon ta fatigue ;
un prompt figé dans une commande s'applique toujours pareil, avec les mêmes
exigences. Tu transformes une bonne formulation, trouvée un jour de lucidité,
en un actif réutilisable par toi et par ton équipe.}

\paragraphe{Concrètement, un skill est un dossier contenant un fichier
\code{SKILL.md}. Le nom du dossier devient le nom de la commande ; le contenu du
fichier décrit ce que l'agent doit faire. Les anciens fichiers
\code{.claude/commands/*.md} restent reconnus, mais les skills sont plus
flexibles : ils peuvent embarquer des exemples, des scripts et des ressources de
référence.}

\section{Créer ton propre skill pour une tâche récurrente}

\paragraphe{Prenons une tâche que tu fais sans cesse : préparer un message de
commit propre à partir des changements en cours. Tu crées un dossier
\code{.claude/skills/commit/}, puis un fichier \code{SKILL.md} dedans.}

\begin{extraitcode}{text}
---
name: commit
description: Préparer un message de commit clair à partir du diff indexé
---

Lis le diff des fichiers indexés (git diff --cached).
Rédige un message de commit en français, format Conventional Commits.
Première ligne : 72 caractères maximum, à l'impératif.
Corps : explique le POURQUOI du changement, pas le QUOI.
N'invente aucun ticket ni numéro. Ne commit pas toi-même : propose seulement.
\end{extraitcode}

\paragraphe{Désormais, \code{/commit} charge ce skill et applique ces consignes.
Tu
remarques le dernier garde-fou : « ne commit pas toi-même ». Une bonne commande
borne son propre périmètre. La commande la plus utile est souvent celle qui dit
aussi à l'agent ce qu'il ne doit \emphase{pas} faire.}

\begin{astuce}
Versionne tes skills dans le dépôt, sous \code{.claude/}. Ils deviennent
alors un bien commun de l'équipe : la même \code{/relecture} s'applique pour
tout le monde, et une amélioration profite à tous au prochain \code{git pull}.
Garde-les courts et mono-tâche ; un skill qui veut tout faire ne fait rien
de fiable.
\end{astuce}

\section{Hooks : déclencher une action à un moment précis du cycle}

\paragraphe{Un skill ou une commande slash, c'est toi qui le déclenches. Un \terme{hook},
lui, se déclenche \emphase{tout seul}, à un moment précis du cycle de l'agent.
C'est une action automatique attachée à un événement : avant qu'une commande
shell s'exécute, après qu'un fichier a été modifié, à la fin d'une réponse.
L'agent n'a rien à décider ; le hook se déclenche, point.}

\paragraphe{La différence est fondamentale. Une commande dépend de ta
discipline : si tu oublies de l'appeler, rien ne se passe. Un hook ne dépend
de personne. C'est ce qui en fait un excellent support de garde-fou :
une règle qui ne s'applique que « quand on y pense » n'est pas une règle.}

\paragraphe{Un hook se configure dans les réglages, en associant un événement à
un ou plusieurs gestionnaires. Voici un modèle réaliste : après une modification
de fichier, Claude Code appelle un script de formatage. Le script reçoit le
contexte de l'action en JSON sur son entrée standard ; il ne dépend donc pas de
variables inventées.}

\begin{extraitcode}{json}
{
  "hooks": {
    "PostToolUse": [
      {
        "matcher": "Write|Edit|MultiEdit",
        "hooks": [
          {
            "type": "command",
            "command": "${CLAUDE_PROJECT_DIR}/.claude/hooks/format-edited-file.sh"
          }
        ]
      }
    ]
  }
}
\end{extraitcode}

\paragraphe{Le script \code{format-edited-file.sh} lit alors le chemin modifié
dans le JSON reçu, par exemple avec \code{jq}, puis lance le formateur du projet.
Les noms exacts des événements et des champs restent à vérifier dans la doc de
ta version, mais la forme importante est celle-ci : configuration dans
\code{settings.json}, script versionné, entrée structurée sur \code{stdin}.}

\section{Cas d'usage concrets : formatage, tests, garde-fous}

\paragraphe{Trois familles de hooks couvrent l'essentiel des besoins réels.}

\begin{liste}
  \element{\important{Formatage automatique.} Après chaque modification de
  fichier, lancer le formateur du projet (Prettier, Black, gofmt). L'agent ne
  produit plus jamais de diff mal indenté : le style est garanti par la machine,
  pas négocié à chaque relecture.}
  \element{\important{Tests à la volée.} Après modification d'un fichier source,
  lancer la suite de tests rapides du module touché. Si un test casse, l'agent
  le voit immédiatement dans le résultat et corrige dans la foulée, au lieu de
  te livrer du code rouge.}
  \element{\important{Garde-fou de sécurité.} Avant qu'une commande shell
  s'exécute, inspecter ce qu'elle contient et bloquer les motifs dangereux
  (\code{rm -rf}, écriture hors du dépôt, \code{git push --force}). Le hook
  refuse l'action avant qu'elle ne parte.}
\end{liste}

\paragraphe{Ce troisième cas relie ce chapitre à ce que tu sais déjà des
permissions. Une permission décide si une action passe ; un hook de
blocage, lui, peut inspecter le \emphase{contenu} de l'action et la refuser
selon une logique fine que tu écris toi-même. Les deux mécanismes se cumulent.}

\begin{extraitcode}{bash}
#!/usr/bin/env bash
# .claude/hooks/block-rm.sh
# Hook PreToolUse : refuse toute suppression récursive forcée.
COMMAND=$(jq -r '.tool_input.command // ""')

if echo "$COMMAND" | grep -qE 'rm +-[a-z]*r[a-z]*f'; then
  jq -n '{
    hookSpecificOutput: {
      hookEventName: "PreToolUse",
      permissionDecision: "deny",
      permissionDecisionReason: "Suppression récursive bloquée par le garde-fou."
    }
  }'
else
  exit 0
fi
\end{extraitcode}

\section{Composer commandes et hooks dans un workflow}

\paragraphe{La vraie puissance vient de la combinaison. Commandes et hooks ne
sont pas deux gadgets isolés : ce sont les deux moitiés d'un même geste
industrialisé. La commande porte ton intention ; les hooks garantissent les
invariants pendant que l'intention s'exécute.}

\paragraphe{Imagine un flux « corriger un bug » de bout en bout. Tu appelles
ta commande \code{/corrige-bug} qui contient le prompt : reproduire, écrire un
test qui échoue, corriger, vérifier. Pendant que l'agent travaille, un hook de
formatage nettoie chaque fichier touché, un hook de test relance la suite après
chaque écriture, et un hook de sécurité veille sur les commandes shell. Tu as
décrit l'intention une fois ; les garde-fous tournent en arrière-plan sans que
tu lèves le petit doigt.}

\paragraphe{Le résultat se relit comme du travail propre : un diff
correctement formaté, des tests verts, aucune commande risquée passée. Tu n'as
pas eu à vérifier chacun de ces points à la main, parce que tu les as outillés
une bonne fois.}

\section{Limites et pièges de l'automatisation prématurée}

\paragraphe{Tout ceci est séduisant, et c'est là le danger. Le piège numéro un
est l'automatisation \emphase{prématurée} : créer une commande pour un geste
que tu n'as fait que deux fois, ou empiler des hooks « au cas où ». Chaque
automatisation est un actif à maintenir. Une commande oubliée qui contient un
vieux prompt périmé fait pire que pas de commande du tout.}

\paragraphe{La règle pragmatique : automatise un geste seulement quand tu l'as
répété assez pour en connaître la bonne forme. Avant ça, tu figerais une
mauvaise version. Laisse le besoin émerger, puis cristallise-le.}

\begin{avertissement}
Un hook qui échoue silencieusement est pire qu'un bug visible. Si ton hook de
test plante pour une raison technique (dépendance manquante, chemin cassé),
l'agent peut interpréter l'erreur de travers et partir corriger un faux
problème. Fais échouer tes hooks \emphase{bruyamment}, avec un message clair sur
la sortie d'erreur, et teste-les comme tu testerais n'importe quel script.
\end{avertissement}

\paragraphe{Deuxième piège : le hook trop lent. Un formateur ou une suite de
tests lancés après chaque écriture, sur un gros projet, peuvent transformer
chaque action de l'agent en attente pénible, et faire exploser la facture vue
au chapitre précédent. Cible tes hooks (un seul module, pas tout le dépôt) et
mesure leur durée.}

\resumeChapitre{
  \element{\important{Skills et commandes} : encapsule un prompt récurrent dans un \code{SKILL.md} versionné — la constance d'une bonne formulation figée vaut plus que la variabilité d'un prompt réécrit de tête.}
  \element{\important{Hooks} : attachés à un événement du cycle, ils s'appliquent sans que tu y penses — formatage après chaque modification, tests après chaque écriture, blocage avant une commande risquée.}
  \element{\important{Composition} : la commande porte l'intention, les hooks garantissent les invariants pendant l'exécution — les deux se cumulent pour un workflow fiable de bout en bout.}
  \element{\important{Automatiser au bon moment} : fige un geste seulement quand tu l'as répété assez pour en connaître la bonne forme — cristalliser trop tôt fige une mauvaise version.}
  \element{\important{Hooks bruyants} : un hook qui échoue silencieusement est pire qu'un bug visible — fais-les échouer avec un message clair et teste-les comme du vrai code.}
}

\paragraphe{Enfin, garde la main. L'automatisation déplace l'effort, elle ne le
supprime pas : tu ne relis plus l'indentation, mais tu dois relire la logique de
tes hooks et de tes commandes. C'est du code, avec ses propres bugs. Bien
outillé, tu as transformé des gestes manuels et faillibles en raccourcis
fiables. L'étape suivante consiste à ouvrir l'agent sur le monde extérieur :
au prochain chapitre, on lui branche de vrais outils externes — bases de
données, gestionnaires de tickets, documentation — grâce au protocole MCP.}
